\documentclass[aspectratio=1610,onlymath]{beamer}
% \documentclass[aspectratio=1610,onlymath,handout]{beamer}

\input{macros-lecture}

\defineTitle{4}{Resolution und Horn-Logik}{23. April 2026}

\begin{document}

\maketitle

\begin{frame}\frametitle{Logik: Glossar}

\begin{tabular}{@{}rll@{}}
\rowcolor{darkblue!10}
Atom & kleinste mögliche Formel & $p\in\Slang{P}$\\
%
Teilformel & Unterausdruck, der Formel ist & $\textsf{Sub}(F)$ \\
%
\rowcolor{darkblue!10}
Wertzuweisung & Funktion & $w:\Slang{P}\to\{\mytrue,\myfalse\}$\\
%
Modell & erfüllende Wertzuweisung & $w\models F$ / $w(F)=\mytrue$\\
%
\rowcolor{darkblue!10}
Literal & Atom oder negiertes Atom & $p$ oder $\neg p$\\
%
Klausel & Disjunktion von Literalen & $L_1\vee\ldots\vee L_n$ \\
%
\rowcolor{darkblue!10}
Monom & Konjunktion von Literalen & $L_1\wedge\ldots\wedge L_n$ \\
%
logische & Teilmengenbeziehung der  & $F\models G$ oder \\[-0.8ex]
Konsequenz & Modellmengen; entspricht $\to$ &  $\mathcal{F}\models G$\\
%
\rowcolor{darkblue!10}
semantische & Gleichheit der  & $F\equiv G$ oder \\[-0.8ex]
\rowcolor{darkblue!10}
Äquivalenz & Modellmengen; entspricht $\leftrightarrow$ &  $\mathcal{F}\equiv\mathcal{G}$\\
%
Tautologie & allgemeingültige Formel & $\models F$

\end{tabular}

\end{frame}


\begin{frame}\frametitle{Rückblick: Vereinfachungen und Normalformen}

Jede Formel ist semantisch äquivalent zu einer, die nur folgende Junktoren verwendet:
\begin{itemize}
\item $\wedge$ und $\neg$
\item $\vee$ und $\neg$
\item und verschiedene andere Kombinationen \ldots\ (z.B. $\to$ und $\neg$; oder auch $\uparrow$)
\end{itemize}\medskip

\emph{Normalformen}
\begin{itemize}
\item \alert{Negationsnormalform}\\
	stets linear
\item \alert{Konjunktive Normalform}\\
	im schlimmsten Fall exponentiell; mit Hilfsatomen linear
\item \alert{Disjunktive Normalform}\\
	im schlimmsten Fall exponentiell
\end{itemize}

\end{frame}



% \begin{frame}\frametitle{Rückblick: Resolution}
% 
% \defbox{
% Gegeben seien zwei Klauseln $K_1$ und $K_2$ für die es ein Atom $p\in\Slang{P}$ gibt mit $p\in K_1$ und $\neg p\in K_2$.
% Die \redalert{Resolvente von $K_1$ und $K_2$ bezüglich $p$} ist die Klausel \alert{$(K_1\setminus\{p\})\cup(K_2\setminus\{\neg p\})$}.
% }\medskip
% 
% \codebox{\emph{Resolution}\\
% Gegeben: Formel $\mathcal{F}$ in Klauselform\\
% Gesucht: Ist $\mathcal{F}$ erfüllbar oder unerfüllbar?\\[1ex]
% 
% \begin{enumerate}[(1)]
% \item Finde ein Klauselpaar $K_1,K_2\in \mathcal{F}$ mit Resolvente $R\notin \mathcal{F}$
% \item Setze $\mathcal{F}\defeq \mathcal{F}\cup\{R\}$
% \item Wiederhole Schritt (1) und (2) bis keine neuen Resolventen gefunden werden können
% \item Falls $\bot\in \mathcal{F}$, dann gib "`unerfüllbar"' aus;\\andernfalls gib "`erfüllbar"' aus
% \end{enumerate}
% }
% \end{frame}

\sectionSlide{Resolution}

\newcommand{\Aname}{Anna}
\newcommand{\Bname}{Barbara}
\newcommand{\Cname}{Chris}

\begin{frame}\frametitle{Resolution}

\emph{Ziel:} Verfahren, um \redalert{(Un)Erfüllbarkeit} einer \redalert{KNF} zu bestimmen
\bigskip

Leider ist (Un)Erfüllbarkeit nicht an einzelnen Klauseln erkennbar.
\medskip

\examplebox{\emph{Beispiel:} Die Logelei aus Vorlesung 1 haben wir mit den Formeln $\mathcal{F}=\{A\leftrightarrow \neg B, B\leftrightarrow \neg C,C\leftrightarrow (\neg A\wedge \neg B)\}$ dargestellt. Wir wollen beweisen, dass
\Bname{} die Wahrheit sagt, also $\mathcal{F}\models B$. Dazu kann man zeigen, dass $\mathcal{F}\cup\{\neg B\}$ unerfüllbar
ist. Bei Anwendung des Distributivgesetztes erhalten wir die folgende KNF:
\begin{align*}
(\neg A\vee \neg B)\wedge (A\vee B)&\wedge{}\\
(\neg B\vee \neg C)\wedge (B\vee C)&\wedge{}\\
(\neg C\vee \neg A)\wedge(A\vee B\vee C)&\wedge{}\\
\neg B
\end{align*}
}

\end{frame}

\begin{frame}\frametitle{Die Klauselform}

Wenn man grundsätzlich mit KNF arbeitet, dann bieten sich einige syntaktische Vereinfachungen an:
\begin{itemize}
\item Eine Klausel $L_1\vee\ldots\vee L_n$ wir dargestellt als Menge $\{L_1,\ldots,L_n\}$
\item Eine Konjunktion von Klauseln $K_1\wedge\ldots\wedge K_\ell$ wird dargestellt als Menge $\{K_1,\ldots,K_\ell\}$
\end{itemize}
Eine Menge von Mengen von Literalen unter dieser Interpretation heißt \redalert{Klauselform}.
\medskip\pause

\examplebox{\emph{Beispiel:} Die KNF\vspace{-3ex}
\begin{align*}
(\neg A\vee \neg B)\wedge (A\vee B)&\wedge{}\\
(\neg B\vee \neg C)\wedge (B\vee C)&\wedge{}\\
(\neg C\vee \neg A)\wedge(A\vee B\vee C)&\wedge{}\\
\neg B
\end{align*}\vspace{-5ex}

kann in Klauselform dargestellt werden als:\vspace{-1ex}
\[\{\{\neg A, \neg B\}, \{A, B\},\{\neg B, \neg C\},\{B, C\},
\{\neg C, \neg A\},\{A, B, C\},\{\neg B\}\}\]
}

\end{frame}

\begin{frame}\frametitle{Ableitungen mittels Resolution}

Beim Resolutionsverfahren leiten wir schrittweise neue Klauseln aus den gegebenen ab.
\bigskip

Ein einzelner Resolutionschritt funktioniert wie folgt:

\defbox{
Gegeben seien zwei Klauseln $K_1$ und $K_2$ für die es ein Atom $p\in\Slang{P}$ gibt mit $p\in K_1$ und $\neg p\in K_2$.
Die \redalert{Resolvente von $K_1$ und $K_2$ bezüglich $p$} ist die Klausel \alert{$(K_1\setminus\{p\})\cup(K_2\setminus\{\neg p\})$}.\medskip

Eine Klausel $R$ ist eine \redalert{Resolvente einer Klauselmenge $\mathcal{K}$} wenn $R$ Resolvente zweier Klauseln  $K_1,K_2\in\mathcal{K}$ ist.
}\pause

\examplebox{\emph{Beispiel:} Resolventen für die Menge $\{\{\neg A, \neg B\},\allowbreak \{A, B\},\allowbreak\{\neg B, \neg C\},\allowbreak\{B, C\},\allowbreak
\{\neg C, \neg A\},\allowbreak\{A, B, C\},\allowbreak\{\neg B\}\}$ sind z.B.
\begin{itemize}
\item $\{B,\neg A\}$ aus den Klauseln $\{B, C\}$ und $\{\neg C,\neg A\}$
\item $\{A\}$ aus den Klauseln $\{A, B\}$ und $\{\neg B\}$
\item $\{B,\neg B\}$ aus den Klauseln $\{B,C\}$ und $\{\neg B, \neg C\}$
\end{itemize}

}

\end{frame}

\begin{frame}\frametitle{Bedeutung von Resolutionsschritten}

Wir betrachten Klauseln $\{L_1,\ldots,L_n,p\}$ und $\{\neg p,M_1,\ldots,M_\ell\}$:\pause
\begin{itemize}
\item $\{L_1,\ldots,L_n,p\}\equiv (L_1\vee\ldots\vee L_n\vee p) \equiv \alert{(\neg L_1\wedge\ldots\wedge\neg L_n)\to p}$\pause
\item $\{\neg p,M_1,\ldots,M_\ell\}\equiv (\neg p\vee M_1\vee\ldots\vee M_\ell)\equiv \alert{p\to (M_1\vee\ldots\vee M_\ell)}$\pause
\end{itemize}\medskip

Daraus folgt unmittelbar \alert{$(\neg L_1\wedge\ldots\wedge\neg L_n)\to (M_1\vee\ldots\vee M_\ell)$}\\
$\leadsto$ dies entspricht der Klausel $\{L_1,\ldots,L_n,M_1,\ldots,M_\ell\}$\pause
\bigskip

\theobox{\emph{Satz:} Wenn $R$ Resolvente der Klauseln $K_1$ und $K_2$ ist, dann gilt $\{K_1,K_2\}\models R$.
}

Resolutionsschritte produzieren also logische Schlüsse. Diese Eigenschaft einer Ableitungsregel wird \redalert{Korrektheit} genannt ("`jede Ableitung ist tatsächlich eine logische Konsequenz"').

\end{frame}

\begin{frame}\frametitle{Die leere Klausel}

Die Resolvente kann eine leere Menge $\emptyset$ sein.\\[1ex]

\examplebox{\emph{Beispiel:} Die Resolvente der Klauseln $\{p\}$ und $\{\neg p\}$ ist leer.
}\medskip

\emph{Was bedeutet so eine \redalert{leere Klausel}?}\pause
\begin{itemize}
\item Klauseln sind Disjunktionen
\item Die leere Klausel entspricht einer Disjunktion von $0$ Literalen
\item Dieser Ausdruck sollte das neutrale Element der Disjunktion sein, d.h. immer den Wert $\myfalse$ annehmen
\end{itemize}
$\leadsto$ Wir bezeichnen die leere Klausel mit $\bot$.
\bigskip\pause

\emph{Was bedeutet es, wenn $\bot$ in Klauselform auftaucht?}\pause
\begin{itemize}
\item Die Klauselform beschreibt eine Konjunktion von Klauseln
\item Wenn eine der Klauseln $\bot$ ist, dann ist die gesamte Formel unerfüllbar
\end{itemize}
$\leadsto$ Die Ableitung von $\bot$ zeigt Unerfüllbarkeit.

\end{frame}

\begin{frame}\frametitle{Das Resolutionskalkül}

Wir können damit das gesamte Verfahren angeben:

\codebox{\emph{Resolution}\\
Gegeben: Formel $\mathcal{F}$ in Klauselform\\
Gesucht: Ist $\mathcal{F}$ erfüllbar oder unerfüllbar?\\[1ex]

\begin{enumerate}[(1)]
\item Finde ein Klauselpaar $K_1,K_2\in \mathcal{F}$ mit Resolvente $R\notin \mathcal{F}$
\item Setze $\mathcal{F}\defeq \mathcal{F}\cup\{R\}$
\item Wiederhole Schritt (1) und (2) bis keine neuen Resolventen gefunden werden können
\item Falls $\bot\in \mathcal{F}$, dann gib "`unerfüllbar"' aus;\\andernfalls gib "`erfüllbar"' aus
\end{enumerate}
}\pause

\emph{Beobachtung:} Unerfüllbarkeit steht fest, sobald $\bot$ abgeleitet \ghost{wurde}\\
$\leadsto$ dann kann man das Verfahren frühzeitig abbrechen
\medskip

Erfüllbarkeit kann dagegen erst erkannt werden, wenn alle Resolventen erschöpfend gebildet worden sind

\end{frame}

\begin{frame}\frametitle{Beispiel}

\begin{tabular}{rlll}
(1) & $\{\neg A, \neg B\}$ && \textcolor{devilscss}{\Aname{} oder \Bname{} lügen.}\\
(2) & $\{A, B\}$  && \textcolor{devilscss}{\Aname{} oder \Bname{} sagen die Wahrheit.} \\
(3) & $\{\neg B, \neg C\}$  && \textcolor{devilscss}{\Bname{} oder \Cname{} lügen.} \\
(4) & $\{B, C\}$ && \textcolor{devilscss}{\Bname{} oder \Cname{} sagen die Wahrheit.} \\
(5) & $\{\neg C, \neg A\}$ && \textcolor{devilscss}{\Cname{} oder \Aname{} lügen.} \\
(6) & $\{A, B, C\}$ && \textcolor{devilscss}{Jemand sagt die Wahrheit.}  \\
(7) & $\{\neg B\}$ && \textcolor{devilscss}{\Bname{} lügt.} \\\pause
(8) & $\{C\}$   &  (4)+(7) & \textcolor{devilscss}{\Cname{} sagt die Wahrheit.} \\\pause
(9) & $\{\neg A\}$ & (8)+(5) & \textcolor{devilscss}{\Aname{} lügt.}\\\pause
(10) & $\{B\}$  & (2)+(9) & \textcolor{devilscss}{\Bname{} sagt die Wahrheit.} \\\pause
(11) & $\bot$ & (10)+(7) & \textcolor{devilscss}{Widerspruch.} 
\end{tabular}\bigskip

$\leadsto$ Es gilt $\{\{\neg A, \neg B\},\{A, B\},\{\neg B, \neg C\},\{B, C\},\{\neg C, \neg A\},\{A, B, C\}\}\models B$

\end{frame}

\begin{frame}[t]\frametitle{Resolution: Eigenschaften}

Resolution eignet sich tatsächlich zum logischen Schließen:\medskip

\theobox{\emph{Satz:} Das Resolutionskalkül hat die folgenden Eigenschaften:
\begin{itemize}
\item \alert{Korrektheit:} Wenn $\bot$ aus $\mathcal{F}$ abgeleitet wird, dann gilt $\mathcal{F}\models\bot$.
\item \alert{Vollständigkeit:} Wenn $\mathcal{F}\models\bot$, dann kann $\bot$ aus $\mathcal{F}$ abgeleitet werden.
\item \alert{Terminierung:} Das Verfahren endet nach endlich vielen Schritten.
\end{itemize}
}\pause

\emph{Beweis:} \alert{Korrektheit} haben wir schon für einzelne Resolutionsschritte gezeigt.
Die Behauptung gilt daher auch, wenn beliebig viele Schritte nacheinander ausgeführt werden (Induktion
mit Hypothese: "`Die neu abgeleitete Klauselmenge folgt aus der ursprünglichen."').

\end{frame}

\begin{frame}[t]\frametitle{Resolution: Eigenschaften (2)}

Resolution eignet sich tatsächlich zum logischen Schließen:\medskip

\theobox{\emph{Satz:} Das Resolutionskalkül hat die folgenden Eigenschaften:
\begin{itemize}
\item \alert{Korrektheit:} Wenn $\bot$ aus $\mathcal{F}$ abgeleitet wird, dann gilt $\mathcal{F}\models\bot$.
\item \alert{Vollständigkeit:} Wenn $\mathcal{F}\models\bot$, dann kann $\bot$ aus $\mathcal{F}$ abgeleitet werden.
\item \alert{Terminierung:} Das Verfahren endet nach endlich vielen Schritten.
\end{itemize}
}

\emph{Beweis:} \alert{Terminierung} ist leicht zu sehen.\pause{}
Jede neu abgeleitete Klausel enthält nur Literale,
die auch in der ursprünglichen Klauselmenge vorkamen. Die Zahl dieser Literale ist endlich; wir bezeichnen sie mit
$\ell$. Es gibt nur $2^\ell$ Klauseln mit diesen Literalen. Man kann also weniger als $2^\ell$ Resolventen hinzufügen bevor das Verfahren terminiert.

\end{frame}

\begin{frame}[t]\frametitle{Resolution: Eigenschaften (3)}

Resolution eignet sich tatsächlich zum logischen Schließen:\medskip

\theobox{\emph{Satz:} Das Resolutionskalkül hat die folgenden Eigenschaften:
\begin{itemize}
\item \alert{Korrektheit:} Wenn $\bot$ aus $\mathcal{F}$ abgeleitet wird, dann gilt $\mathcal{F}\models\bot$.
\item \alert{Vollständigkeit:} Wenn $\mathcal{F}\models\bot$, dann kann $\bot$ aus $\mathcal{F}$ abgeleitet werden.
\item \alert{Terminierung:} Das Verfahren endet nach endlich vielen Schritten.
\end{itemize}
}

\emph{Beweis:} \alert{Vollständigkeit} ist etwas aufwändiger.\medskip

Vorgehen: Wir zeigen, wie
man eine erfüllende Zuweisung finden kann, wenn die leere Klausel nicht abgeleitet wurde.

\end{frame}

\begin{frame}\frametitle{Widerspruchsfreie Resolution $\leadsto$ Modell}

% Wenn $\bot$ nicht abgeleitet wurde, dann kann man ein Modell der Formel wie folgt finden.
% \medskip

Für jedes Atom $p$, das in der Formel vorkommt, bestimmen wir einen Wahrheitswert $w(p)$:\medskip

\codebox{
Gegeben: widerspruchsfreie Klauselmenge $\mathcal{F}$ nach Resolution\\[1ex]

Für jedes Atom $p$ in $\mathcal{F}$ (in beliebiger Reihenfolge):
\begin{enumerate}[(a)]
\item Wenn $\{p\}\in\mathcal{F}$, definiere $w(p)\defeq\mytrue$
\item Wenn $\{\neg p\}\in\mathcal{F}$, definiere $w(p)\defeq\myfalse$
\item Wenn weder $\{p\}\in\mathcal{F}$ noch $\{\neg p\}\in\mathcal{F}$, dann setze $w(p)\defeq\mytrue$
und $\mathcal{F}\defeq\mathcal{F}\cup\{\{p\}\}$ und führe nochmals Resolution durch, bis keine weiteren Klauseln abgeleitet
werden.
\end{enumerate}
}\smallskip

Die so bestimmte Wertzuweisung $w$ ist ein Modell für $\mathcal{F}$.\bigskip

{\tiny  Beweisskizze:
\begin{enumerate}[(1)]
\item In Schritt (c) wird nie die leere Klausel abgeleitet (Intuition: Dazu wäre eine Klausel $\{\neg p\}$ nötig, die es aber in diesem Fall nicht geben darf und die auch nicht neu abgeleitet werden kann $\leadsto$ Induktion)
\item Dank (1) enthält $\mathcal{F}$ auch nach Bestimmung aller Werte keine leere Klausel (einfache Induktion)
\item Daher werden alle Klauseln in $\mathcal{F}$ erfüllt (Wenn eine nicht erfüllt wäre, dann könnte man aus ihr unter Verwendung der einelementigen Klauseln $\bot$ ableiten, was aber laut (2) nicht abgeleitet wird)
\end{enumerate}}
% 
% Man kann zeigen: $w\models\mathcal{F}$, insbesondere ist $w$ ein Modell der ursprünglichen Formel

% Wir benötigen eine Wertzuweisung für jedes Atom, das in der Formel vorkommt.
% 
% \begin{itemize}
% \item Wenn eine Klausel $\{p\}$ oder $\{\neg p\}$ abgeleitet wurde, dann kann man die
% Wertzuweisung ablesen ($p\mapsto\mytrue$ bzw. $p\mapsto\myfalse$)
% \item Wenn für ein Atom $p$ weder $\{p\}$ noch $\{\neg p\}$ abgeleitet wurde, dann können wir eines
% davon frei wählen (z.B. $\{p\}$) und mit dieser zusätzlichen Klausel Resolution durchführen (dies kann
% keinen Widerspruch erzeugen!)
% \end{itemize}

\end{frame}


\begin{frame}[t]\frametitle{Resolution: Eigenschaften (4)}

Resolution eignet sich tatsächlich zum logischen Schließen:\medskip

\theobox{\emph{Satz:} Das Resolutionskalkül hat die folgenden Eigenschaften:
\begin{itemize}
\item \alert{Korrektheit:} Wenn $\bot$ aus $\mathcal{F}$ abgeleitet wird, dann gilt $\mathcal{F}\models\bot$.
\item \alert{Vollständigkeit:} Wenn $\mathcal{F}\models\bot$, dann kann $\bot$ aus $\mathcal{F}$ abgeleitet werden.
\item \alert{Terminierung:} Das Verfahren endet nach endlich vielen Schritten.
\end{itemize}
}

\emph{Beweis:} \alert{Vollständigkeit} ist etwas aufwändiger.\medskip

Vorgehen: Wir zeigen, wie
man eine erfüllende Zuweisung finden kann, wenn die leere Klausel nicht abgeleitet wurde.\qed
\bigskip

\emph{Anmerkung:} "`Vollständigkeit"' bezieht sich nur auf die Ableitung von $\bot$. Resolution kann nicht
jede beliebige Klausel ableiten, die logisch folgt!

\end{frame}

\begin{frame}\frametitle{Beispiel: widerspruchsfreie Resolution}

\begin{minipage}{4cm}
\tiny
\begin{tabular}[t]{rlll}
(1) & $\{\neg A, \neg B\}$ &\\
(2) & $\{A, B\}$  &\\
(3) & $\{\neg B, \neg C\}$  &\\
(4) & $\{B, C\}$ &\\
(5) & $\{\neg A,\neg C\}$ &\\
(6) & $\{A, B, C\}$ &\\
%
(7) & $\{B,\neg B\}$ & (2)+(1)\\
(8) & $\{B,\neg C\}$ & (2)+(5)\\
(9) & $\{A,\neg A\}$ & (2)+(1)\\
(10) & $\{A,\neg C\}$ & (2)+(3)\\
(11) & $\{\neg A,C\}$ & (4)+(1)\\
(12) & $\{C,\neg C\}$ & (4)+(3)\\
(13) & $\{\neg A,B\}$ & (4)+(5)\\
(14) & $\{B,\neg B,C\}$ & (6)+(1)\\
(15) & $\{B,C,\neg C\}$ & (6)+(5)\\
(16) & $\{A,\neg A, C\}$ & (6)+(1)\\
(17) & $\{A,C, \neg C\}$ & (6)+(3)\\
(18) & $\{A,B,\neg B\}$ & (6)+(3)\\
(19) & $\{A,\neg A, B\}$ & (6)+(3)\\
%
(20) & \textcolor<2->{darkred}{$\{B\}$} & (2)+(13)\\
(21) & \textcolor<2->{darkred}{$\{\neg C\}$} & (8)+(3)\\
(22) & \textcolor<2->{darkred}{$\{\neg A\}$} & (13)+(1)\\
(23) & $\{\neg A,B,\neg B\}$ & (14)+(5)\\
(24) & $\{\neg A,C,\neg C\}$ & (15)+(1)
\end{tabular}\end{minipage}
\begin{minipage}{5cm}{\tiny
\begin{tabular}[t]{rlll}
(25) & $\{A,\neg A,\neg B\}$ & (16)+(3)\\
(26) & $\{\neg B,C,\neg C\}$ & (17)+(1)\\
(27) & $\{B,\neg B,\neg C\}$ & (18)+(5)\\
(28) & $\{A,\neg A,\neg C\}$ & (19)+(3)\\
%
(29) & $\{\neg A,\neg B,\neg C\}$ & (23)+(3)\\
(30) & $\{A,\neg A,B,\neg B\}$ & (6)+(29)\\
(31) & $\{A,\neg A,C,\neg C\}$ & (6)+(29)\\
(32) & $\{B,\neg B,C,\neg C\}$ & (6)+(29)\\
%
(33) & $\{A,\neg A, B,\neg B,C,\neg C\}$ & (30)+(31)\\
(34) & $\{\neg A, B,\neg B,C,\neg C\}$ & (33)+(22)\\
(35) & $\{A, \neg A, B,C,\neg C\}$ & (20)+(33)\\
(36) & $\{A, \neg A, B,\neg B,\neg C\}$ & (33)+(21)\\
(37) & $\{\neg A, B,C,\neg C\}$ & (20)+(34)\\
(38) & $\{\neg A, B,\neg B,\neg C\}$ & (34)+(21)\\
(39) & $\{A, \neg A, B,\neg C\}$ & (35)+(21)\\
(40) & $\{\neg A, B,\neg C\}$ & (39)+(22)\\
(41) & $\{A, B,\neg B,C,\neg C\}$ & (2)+(33)\\
(42) & $\{A,\neg A, \neg B,C,\neg C\}$ & (33)+(3)\\
(43) & $\{A,\neg A, B,\neg B,C\}$ & (4)+(33)\\
(44) & \ldots und so weiter
\end{tabular}}\pause
\vspace{5mm}

Es ergibt sich ein Modell $w$ mit
$w(A)=\myfalse$, $w(B)=\mytrue$ und $w(C)=\myfalse$.
\end{minipage}

\end{frame}

\begin{frame}\frametitle{Resolution: Komplexität}

Leider können \redalert{im schlimmsten Fall exponentiell} viele Resolventen gebildet werden
(wobei die Formel erfüllbar ist, so dass kein frühzeitiger Abbruch möglich ist)
\medskip

\textcolor{devilscss}{(Die Verwendung von Hilfsatomen zur Ermittlung einer linear großen Klauselform löst dieses Problem nicht)}
\bigskip

\begin{itemize}
\item Bezüglich der algorithmischen Komplexität ist Resolution schlimmstenfalls \alert{nicht besser als andere Verfahren}.
\item Das hier vorgestellte einfachste Verfahren ist \alert{hoffnungslos ineffizient}.
\item Praktische Implementierungen verwenden in der Regel \alert{keine (reine) Resolution}, bilden aber manchmal Resolventen
\end{itemize}
\medskip

\anybox{strongyellow}{
Resolution ist dennoch bedeutsam, da sie gut auf Prädikatenlogik und verwandte Logiken anwendbar ist, was für andere aussagenlogische Verfahren nicht zutrifft.\\
(Auch da kommen allerdings viele weitere Optimierungen zum Einsatz.)}

\end{frame}

% \begin{frame}\frametitle{Aussagenlogisches Schließen in der Praxis}
% 
% Systeme zum aussagenlogischen Schließen heißen \redalert{SAT-Solver}.
% \bigskip
% 
% Moderne SAT-Solver \ldots
% \begin{enumerate}[$\ldots$]
% \item sind \alert{stark optimiert} und können Probleme mit Millionen von Atomen lösen
% \item verwenden in der Regel \alert{keine (reine) Resolution}, bilden aber manchmal Resolventen
% \item werden auch \alert{jenseits der Aussagenlogik} verwendet, indem Probleme zunächst in Aussagenlogik übersetzt werden
% \end{enumerate}\bigskip
% 
% \end{frame}




\sectionSlide{Horn-Logik}

\begin{frame}\frametitle{Horn-Logik}

Bisher waren alle unsere Ansätze für aussagenlogisches Schließen exponentiell -- geht es auch einfacher?
\bigskip\pause

\emph{Idee:} Beschränke die Form von Formeln
\bigskip

\defbox{Eine \redalert{Horn-Klausel}${}^*$ ist eine Klausel, die höchstens ein nichtnegiertes Literal enthält.
Eine \redalert{Horn-Formel} ist eine Formel in KNF, welche nur Horn-Klauseln enthält.\\
{\tiny ${}^*$) nach Alfred Horn, 1918--2001, US-amerikanischer Mathematiker}}

\pause
\examplebox{\emph{Beispiele:}
\begin{itemize}
\item $\neg p\vee \neg q \vee r$ ist eine Horn-Klausel
\item $q\vee p$ ist keine Horn-Klausel
\item $p$ ist eine Horn-Klausel
\item $\neg p\vee \neg q$ ist eine Horn-Klausel
\item $p\vee p$ ist keine Horn-Klausel (aber äquivalent zu einer)
\end{itemize}
}

\end{frame}

\begin{frame}\frametitle{Horn-Klauseln als Implikationen}

Jede Horn-Klausel kann als eine Implikation ohne Negation und Disjunktion ausgedrückt werden
\bigskip

\examplebox{\emph{Beispiele:}
\begin{align*}
\neg p\vee \neg q \vee r \qquad \equiv &&(p\wedge q) &\to r\\
p \qquad \equiv &&\top&\to p\\
\neg p\vee \neg q \qquad \equiv &&(p\wedge q)&\to \bot
\end{align*}
}

\begin{itemize}
\item Wir verwenden $\top$ für die leere Prämisse und $\bot$ für die leere Konsequenz (weil so die gewünschte logische Äquivalenz gilt)
\item Als Implikationen geschriebene Horn-Klauseln werden oft als \redalert{(Horn-)Regeln} bezeichnet
\item Es ist üblich, die Klammern der Konjunktion in der Prämisse wegzulassen
\end{itemize}

\end{frame}

% \begin{frame}\frametitle{Ausblick}
% 
% \emph{Nächste Vorlesung:} Wir werden logisches Schließen für Horn-Regeln noch etwas genauer betrachten
% \begin{itemize}
% \item Resolution kann in diesem Fall spezialisiert werden
% \item Dadurch erhält man polynomielle Algorithmen
% \end{itemize}
% 
% \end{frame}


% \begin{frame}\frametitle{Rückblick: Horn-Logik}
% 
% \defbox{Eine \redalert{Horn-Klausel}${}^*$ ist eine Klausel, die höchstens ein nichtnegiertes Literal enthält.
% Eine \redalert{Horn-Formel} ist eine Formel in KNF, welche nur Horn-Klauseln enthält.\\
% {\tiny ${}^*$) nach Alfred Horn, 1918--2001, US-amerikanischer Mathematiker}}
% 
% Man kann Horn-Formeln als Implikationen auffassen \\
% (sogenannte \redalert{Horn-Regeln}):
% 
% \examplebox{Beispiele:\vspace{-4ex}
% \begin{align*}
% \neg p\vee \neg q \vee r \qquad \equiv &&(p\wedge q) &\to r\\
% p \qquad \equiv &&\top&\to p\\
% \neg p\vee \neg q \qquad \equiv &&(p\wedge q)&\to \bot
% \end{align*}
% }
% 
% \begin{itemize}
% \item Wir verwenden $\top$ für die leere Prämisse und $\bot$ für die leere Konsequenz (weil so die gewünschte logische Äquivalenz gilt)
% % \item Als Implikationen geschriebene Horn-Klauseln werden oft als \redalert{(Horn-)Regeln} bezeichnet
% \item Es ist üblich, die Klammern der Konjunktion in der Prämisse wegzulassen
% \end{itemize}
% 
% \end{frame}

\begin{frame}\frametitle{Beispiel}

Viele einfache rekursive Algorithmen können in Horn-Aussagenlogik beschrieben werden.
\bigskip

\examplebox{\emph{Beispiel:} In der Lehrveranstaltung \href{https://iccl.inf.tu-dresden.de/web/FS2025}{\alert{Formale Systeme}} berechneten wir die Variablen $V_\epsilon$ einer CFG, welche in $\epsilon$ umgeschrieben werden können (siehe Vorlesung 2, Folie 31). Wir betrachten die folgende Menge $\mathcal{F}$ von Horn-Regeln:
\begin{align*}
\top & \to p_{\Snterm{A}} & \text{für jede Grammatikregel $\Snterm{A}\to\epsilon$}\\
p_{\Sntermsub{A}{1}}\wedge\ldots\wedge p_{\Sntermsub{A}{n}} & \to p_{\Snterm{B}} & \text{für jede Grammatikregel $\Snterm{B}\to\Sntermsub{A}{1}\cdots\Sntermsub{A}{n}$}
\end{align*}
Dann gilt $\mathcal{F}\models p_{\Snterm{A}}$ genau dann wenn $\Snterm{A}\in V_\epsilon$.
}

\end{frame}


\begin{frame}\frametitle{Resolution bei Horn-Klauseln}

Die Resolution von Horn-Klauseln = Verknüpfung von Regeln:
%
\[\frac{p_1\wedge\ldots\wedge p_n\to \redalert{q}\qquad\redalert{q}\wedge q_1\wedge\ldots\wedge q_m\to r}{p_1\wedge\ldots\wedge p_n\wedge q_1\wedge\ldots\wedge q_m\to r}\]\pause

Man kann zeigen (ohne Beweis):

\theobox{\emph{Satz:} Das Resolutionskalkül für Horn-Formeln bleibt auch dann vollständig, wenn man sich auf Resolventen
beschränkt, bei denen eine Klausel die Form $\{p\}$ (also $\top\to p$) hat.
}

Diese Einschränkung entspricht der Entfernung von bereits erfüllten Vorbedingungen aus der Prämisse einer Regel:
\[\frac{\top\to\redalert{q}\qquad\redalert{q}\wedge q_1\wedge\ldots\wedge q_m\to r}{q_1\wedge\ldots\wedge q_m\to r}\]

\end{frame}

\begin{frame}\frametitle{Hyperresolution = Anwendung von Regeln}

\emph{Es ist leicht zu sehen:} Im eingeschränkten Resolutionskalkül für Horn-Formeln kann eine Resolvente aus einer Regel $q_1\wedge\ldots\wedge q_m\to r$
nur dann zur Ableitung von $\bot$ nützlich sein, wenn letztlich alle Vorbedingungen $q_1,\ldots, q_m$ erfüllt sind.

$\leadsto$ Wir können die Resolution aufschieben, bis dies gegeben ist
\bigskip\pause

\defbox{%
\redalert{Hyperresolution:} Mehrere Klauseln mit positiven Literalen werden parallel mit einer Klausel mit vielen negativen Literalen resolviert.
%
\[\frac{\top\to q_1\qquad\ldots\qquad\top\to q_m\qquad q_1\wedge\ldots\wedge q_m\to r}{\top\to r}\]}
%
Dies entspricht der intuitiven Idee einer \alert{Regelanwendung}: "`Wenn die Regel $q_1\wedge\ldots\wedge q_m\to r$ und die Atome $q_1$, \ldots, $q_m$ gelten, dann kann $r$ abgeleitet werden."'

\end{frame}

\begin{frame}\frametitle{Komplexität des Schließens mit Horn-Formeln}

\emph{Vorteil der Hyperresolution bei Horn-Klauseln:}
\begin{itemize}
\item Jede Resolvente hat die Form $\top\to p$
\item Es gibt nur linear viele solche Resolventen (für Atome $p$, die in der Formel vorkommen)
\end{itemize}
$\leadsto$ Resolution muss nach linear vielen Schritten terminieren
\pause\bigskip

Dies führt zu einer polynomiellen Laufzeit, nicht zu einer linearen (da jeder Resolutionsschritt einige Zeit benötigt).
\medskip

Es geht aber auch noch besser:

\theobox{\emph{Satz (Dowling \& Gallier):} Die Erfüllbarkeit einer Horn-Formel kann in linearer Zeit entschieden werden.}

(Ohne Beweis)
% \bigskip

%%% Folgendes ist im Kontext von NP relevant:
% \emph{Intuition:} Horn-Logik ist die \redalert{deterministische} Form der Aussagenlogik (da Regeln immer genau ein
% Atom implizieren). 

\end{frame}

% 
% \begin{frame}\frametitle{Horn-Logik}
% 
% Vereinfachung des logischen Schließens durch Beschränkung auf Horn-Logik:
% \begin{itemize}
% \item (Hyper-)Resolution liefert polynomiellen Algorithmus
% \item Linearer Algorithmus möglich 
% \end{itemize}\bigskip
% 
% 
% \end{frame}



% \begin{frame}\frametitle{Kompaktheit}
% 
% Eine interessante Eigenschaft der Aussagenlogik folgt aus dem, was wir bereits wissen:
% 
% \theobox{\emph{Satz (Kompaktheit der Aussagenlogik):} Wenn $\mathcal{F}$ eine unendliche Formelmenge ist und $\mathcal{F}\models G$ gilt, dann gibt es eine endliche Teilmenge $\mathcal{F}'\subseteq\mathcal{F}$ mit $\mathcal{F}'\models G$.}
% 
% \emph{Beweis:} Die Vollständigkeit der Resolution 
% 
% \end{frame}


\begin{frame}\frametitle{Zusammenfassung und Ausblick}

\redalert{Resolution} ist ein bekanntes Widerlegungskalkül, welches über die Aussagenlogik hinaus von Bedeutung ist
\bigskip

Wie andere Beweisverfahren für Aussagenlogik benötigt Resolution schlimmstenfalls \redalert{exponentiell viel Zeit}%;
% bei \redalert{Horn-Formeln} gibt es dagegen polynomielle Schlussfolgerungsalgorithmen
\bigskip

\redalert{Horn-Logik} erlaubt logisches Schließen in polynomieller Zeit, z.B. mit Resolution.
\bigskip

\anybox{yellow}{
Offene Fragen:
\begin{itemize}
\item Gibt es noch bessere Methoden zum aussagenlogischen Schließen?
\item Was hat Aussagenlogik mit Sprachen, Berechnung und TMs zu tun?
\end{itemize}
}
\end{frame}

\end{document}
